\begin{ouverture}[calcul]{10}{Calculer avec les décimaux}{Partie 2 · Fractions et nombres décimaux}
{Sommes · Produits · Divisions · Rangs}
\objectif{Additionner, soustraire et multiplier des décimaux par un entier.}
\objectif{Multiplier et diviser par 10, 100, 1 000.}
\objectif{Calculer doubles, moitiés et quotients décimaux utiles.}
\prerequis{Opérations posées et rangs jusqu’aux millièmes.}
\end{ouverture}
\pagestyle{otMain}\essentiel\fontsize{11}{15}\selectfont
\phantomsection\label{cours-10-1}
\essSection{1}{Sommes, différences, produits}
\begin{savoirfaire}[Aligner et estimer]Virgules alignées, rangs complétés : $12{,}7+3{,}85=16{,}55$ et $20-4{,}35=15{,}65$. L’estimation encadre : $12{,}7+3{,}85$ est proche de $13+4=17$.\end{savoirfaire}
\begin{exemple}$6{,}38\times7$ : calculer $638\times7=4\,466$, puis replacer deux chiffres après la virgule : $44{,}66$. Estimation : $6{,}4\times7\approx45$.\end{exemple}
\begin{propriete}Multiplier par 10, 100 ou $1\,000$ fait monter les chiffres de 1, 2 ou 3 rangs : $3{,}456\times100=345{,}6$. Diviser fait descendre : $47{,}2:100=0{,}472$.\end{propriete}
\begin{savoirfaire}[Procédures de tête]Doubles et moitiés : le double de $3{,}45$ est $6{,}9$ ; la moitié de $7{,}3$ est $3{,}65$ (la moitié de 73 dixièmes fait 36 dixièmes et demi, soit 365 centièmes). Multiplier par 5 : multiplier par 10 puis prendre la moitié ; par 50 : par 100 puis la moitié.\end{savoirfaire}
\clearpage
\phantomsection\label{cours-10-2}
\essSection{2}{Diviser un décimal}
\begin{exemple}$25{,}5:3$ : 25 unités donnent 8, reste 1 ; 15 dixièmes donnent 5 dixièmes. Quotient : $8{,}5$. Contrôle : $8{,}5\times3=25{,}5$.\end{exemple}
\begin{savoirfaire}[Diviser par 4 ou par 8]Diviser par 4 : deux moitiés successives. $18:4$ : moitié 9, moitié $4{,}5$. Diviser par 8 : trois moitiés. Ces chemins évitent de poser l’opération.\end{savoirfaire}
\begin{propriete}Le quotient peut demander des rangs supplémentaires : $7{,}3:5$ : 7 unités donnent 1 ; 23 dixièmes donnent 4 ; reste 3 dixièmes, soit 30 centièmes ; 30 centièmes donnent 6. Quotient exact : $1{,}46$.\end{propriete}
\begin{attention}Multiplier un nombre par un décimal plus petit que 1 le diminue, et diviser par 10 le rapetisse : les automatismes « multiplier agrandit toujours » cessent d’être vrais hors des entiers. Les rangs, eux, ne mentent jamais.\end{attention}

\clearpage\rubrique{Découverte}
\begin{blocLivre}{ex-10-1}\begin{exercice}[1]{Somme et différence}
Poser et calculer $12{,}7+3{,}85$, puis $20-4{,}35$.
\end{exercice}\end{blocLivre}
\begin{blocLivre}{ex-10-2}\begin{exercice}[1]{Un produit posé}
Poser et calculer $6{,}38\times7$.
\end{exercice}\end{blocLivre}
\begin{blocLivre}{ex-10-3}\begin{exercice}[1]{Par 10, 100, 1 000}
Calculer $3{,}456\times100$ ; $0{,}07\times1\,000$ ; $47{,}2:100$ ; $5:1\,000$.
\end{exercice}\end{blocLivre}
\begin{blocLivre}{ex-10-4}\begin{exercice}[1]{Doubles et moitiés}
Donner le double de $3{,}45$ ; la moitié de $7{,}3$ ; la moitié de $0{,}9$.
\end{exercice}\end{blocLivre}
\begin{blocLivre}{ex-10-5}\begin{exercice}[1]{Multiplier par 5 et par 50}
Calculer de tête $4{,}6\times5$, puis $2{,}4\times50$.
\end{exercice}\end{blocLivre}

\rubrique{Entraînement}
\begin{blocLivre}{ex-10-6}\begin{exercice}[2]{Division en ligne}
Calculer $25{,}5:3$, puis contrôler par la multiplication.
\end{exercice}\end{blocLivre}
\begin{blocLivre}{ex-10-7}\begin{exercice}[2]{Diviser par 4 et par 8}
Calculer $18:4$ et $18:8$ par moitiés successives.
\end{exercice}\end{blocLivre}
\begin{blocLivre}{ex-10-8}\begin{exercice}[2]{Le quotient qui s’allonge}
Calculer le quotient exact de $7{,}3:5$.
\end{exercice}\end{blocLivre}
\begin{blocLivre}{ex-10-9}\begin{exercice}[2]{Compléments et regroupements}
Calculer $4{,}35+2{,}8+0{,}65$ en regroupant astucieusement, puis compléter : $6{,}25+\square=10$.
\end{exercice}\end{blocLivre}
\begin{blocLivre}{ex-10-10}\begin{exercice}[2]{L’estimation qui sauve}
Sans poser le calcul, choisir le bon résultat de $19{,}8\times6$ parmi : $1{,}188$ ; $11{,}88$ ; $118{,}8$ ; $1\,188$. Justifier.
\end{exercice}\end{blocLivre}

\rubrique{Problèmes}
\begin{blocLivre}{ex-10-11}\begin{exercice}[3]{Les courses du club}
Le club achète 6 ballons à $12{,}45$ euros l’un et 3 pompes à $4{,}20$ euros l’une.\begin{questions}
\item Calculer la dépense totale.
\item Le club paie avec 90 euros. Combien lui rend-on ?
\end{questions}
\end{exercice}\end{blocLivre}
\begin{blocLivre}{ex-10-12}\begin{exercice}[3]{Le carrelage du couloir}
Un couloir mesure $7{,}2$ m. On le carrelle avec des dalles de $0{,}4$ m posées bout à bout. Combien de dalles faut-il sur la longueur ?
\end{exercice}\end{blocLivre}
\begin{blocLivre}{ex-10-13}\begin{exercice}[3]{Le partage du week-end}
Trois amis se partagent équitablement $25{,}5$ euros gagnés au vide-grenier, puis chacun dépense $5{,}90$ euros. Combien reste-t-il à chacun ?
\end{exercice}\end{blocLivre}
\begin{blocLivre}{ex-10-14}\begin{exercice}[3]{La moyenne du relais}
Quatre coureurs parcourent ensemble $10{,}8$ km en se partageant également la distance. Quelle distance parcourt chacun ? L’un d’eux ne court que $2{,}5$ km : qui doit compenser, et de combien ?
\end{exercice}\end{blocLivre}
\begin{blocLivre}{ex-10-15}\begin{exercice}[3]{Multiplier rapetisse ?}
Vrai ou faux : « Multiplier un nombre par 10 l’agrandit toujours » ; « Diviser un nombre par 10 l’agrandit parfois ». Justifier avec $0{,}3$.
\end{exercice}\end{blocLivre}

\clearpage
